\section*{Capítulo \ref{ch:g10:cell-metabolism} --- El metabolismo celular}

\begin{solution}{exo:g10:cell-metabolism:1}
El \omterm{def:g10:cell-metabolism:metabolism}{metabolismo} es el conjunto de todas las reacciones químicas de una
\omterm{def:g10:cells-common-unit:cell}{célula}, las que degradan moléculas para obtener energía y las que
construyen las propias. Una \omterm{def:g10:cell-metabolism:metabolism}{enzima} es una \omterm{def:g10:chemistry-of-life:families}{proteína} que acelera una
reacción concreta, uniéndose a su sustrato y liberando el producto, sin
consumirse.
\end{solution}

\begin{solution}{exo:g10:cell-metabolism:2}
Respiración: $\mathrm{C_6H_{12}O_6} + 6\,\mathrm{O_2} \to
6\,\mathrm{CO_2} + 6\,\mathrm{H_2O}$ + energía. \omterm{def:g10:cell-metabolism:autotrophy}{Fotosíntesis}:
$6\,\mathrm{CO_2} + 6\,\mathrm{H_2O}$ + luz $\to \mathrm{C_6H_{12}O_6}
+ 6\,\mathrm{O_2}$. El balance global de materia está invertido; pero
las reacciones, las \omterm{def:g10:cell-metabolism:metabolism}{enzimas}, los lugares (\omterm{def:g10:cells-common-unit:organelle}{mitocondria} frente a
\omterm{def:g10:cells-common-unit:organelle}{cloroplasto}) y las fuentes de energía (enlaces químicos frente a luz)
son completamente distintos.
\end{solution}

\begin{solution}{exo:g10:cell-metabolism:3}
\omterm{def:g10:cell-metabolism:autotrophy}{Heterótrofos}: la seta, la \omterm{def:g10:cells-common-unit:cell}{célula} de raíz de zanahoria, la bacteria del
yogur, la \omterm{def:g10:cells-common-unit:cell}{célula} muscular humana. \omterm{def:g10:cell-metabolism:autotrophy}{Autótrofos}: el musgo, la clorela.
\end{solution}

\begin{solution}{exo:g10:cell-metabolism:4}
El dióxido de carbono infla el globo; el etanol se acumula en el
líquido. Al matraz le falta oxígeno: con aire burbujeando, la levadura
respiraría, no produciría etanol y gastaría mucho menos azúcar.
\end{solution}

\begin{solution}{exo:g10:cell-metabolism:5}
La catalasa es una \omterm{def:g10:chemistry-of-life:families}{proteína}; hervirla destruye su forma y con ella su
actividad, y el agua oxigenada no se descompone deprisa por sí sola.
\end{solution}

\begin{solution}{exo:g10:cell-metabolism:6}
$(7.9 - 0.9)/7 = \qty{1.0}{mg/L}$ por minuto. Se aplana porque el
oxígeno se ha agotado; la levadura pasa entonces a la \omterm{def:g10:cell-metabolism:fermentation}{fermentación}, que
no consume oxígeno.
\end{solution}

\begin{solution}{exo:g10:cell-metabolism:7}
$0.8 + 0.2 = \qty{1.0}{mg/L}$ por minuto: la pendiente a la luz es la
producción menos la respiración que continúa con luz.
\end{solution}

\begin{solution}{exo:g10:cell-metabolism:8}
\qty{18}{g} son \qty{0.1}{mol} de glucosa; necesita \qty{0.6}{mol} de
oxígeno, $0.6 \times 32 = \qty{19.2}{g}$, y desprende \qty{0.6}{mol} de
dióxido de carbono, $0.6 \times 44 = \qty{26.4}{g}$.
\end{solution}

\begin{solution}{exo:g10:cell-metabolism:9}
Los guisantes en germinación respiran sus reservas y liberan calor; los
guisantes hervidos están muertos y sus \omterm{def:g10:cell-metabolism:metabolism}{enzimas} destruidas, así que no
hay reacción ni calor. Cerrado herméticamente, el oxígeno se agotaría y
la respiración se detendría.
\end{solution}

\begin{solution}{exo:g10:cell-metabolism:10}
La mitad expuesta se tiñe de azul muy oscuro (almidón fabricado por
\omterm{def:g10:cell-metabolism:autotrophy}{fotosíntesis} a la luz); la mitad tapada no. Las 48 horas de oscuridad
permiten que la hoja gaste el almidón que ya tenía, de modo que
cualquier almidón encontrado se deba a la luz del día.
\end{solution}

\begin{solution}{exo:g10:cell-metabolism:11}
La \omterm{def:g10:cell-metabolism:fermentation}{fermentación} rinde unas quince veces menos ATP por glucosa que la
respiración; para obtener el ATP que necesita, una \omterm{def:g10:cells-common-unit:cell}{célula} que fermenta
debe degradar quince veces más azúcar por minuto.
\end{solution}

\begin{solution}{exo:g10:cell-metabolism:12}
Durante la carrera, la necesidad de ATP del músculo superó el oxígeno
que la sangre podía entregar; las \omterm{def:g10:cells-common-unit:cell}{células} pasaron a la \omterm{def:g10:cell-metabolism:fermentation}{fermentación}
láctica, cuyo producto, el ácido láctico, se acumuló. En un trote suave
el aporte de oxígeno va al ritmo, la respiración sola proporciona el
ATP y no se produce ácido.
\end{solution}

\begin{solution}{exo:g10:cell-metabolism:13}
Ha cambiado el ambiente (sin luz, con alimento orgánico); los \omterm{def:g10:universal-dna:gene}{genes} no.
Los descendientes reconstruyen \omterm{def:g10:cells-common-unit:organelle}{cloroplastos}, así que las instrucciones
no se perdieron nunca: la pérdida fue un cambio del equipamiento en
uso, no de las posibilidades que permiten los \omterm{def:g10:universal-dna:gene}{genes}.
\end{solution}

\begin{solution}{exo:g10:cell-metabolism:14}
La planta hace la \omterm{def:g10:cell-metabolism:autotrophy}{fotosíntesis} a la luz y aporta oxígeno y \omterm{def:g10:chemistry-of-life:mineral-organic}{materia
orgánica} (hojas que el caracol come); el caracol respira y aporta
dióxido de carbono a la planta. A oscuras la planta solo respira: ambos
consumen oxígeno, nadie lo produce, y los dos mueren en unos días.
\end{solution}

\begin{solution}{exo:g10:cell-metabolism:15}
$30/2 = 15$ veces más glucosa por minuto. Los bodegueros quieren
etanol, que solo produce la \omterm{def:g10:cell-metabolism:fermentation}{fermentación}; una levadura que respirara
convertiría el azúcar en dióxido de carbono, agua y más levadura, y
nada de alcohol.
\end{solution}

\begin{solution}{pb:g10:cell-metabolism:1}
\textbf{1.} \omterm{def:g10:cell-metabolism:fermentation}{Fermentación} alcohólica, porque la cuba cerrada no
contiene oxígeno: $\mathrm{C_6H_{12}O_6} \to 2\,\mathrm{C_2H_5OH} +
2\,\mathrm{CO_2}$.

\textbf{2.} $200/180 \approx \qty{1.11}{mol}$.

\textbf{3.} Etanol $2 \times 1.11 = \qty{2.22}{mol}$, es decir $2.22
\times 46 \approx \qty{102}{g}$; volumen $102/0.79 \approx
\qty{129}{mL}$; alrededor del 13\% de volumen.

\textbf{4.} $2.22 \times 44 \approx \qty{98}{g}$ de dióxido de
carbono; $2.22 \times 24 \approx \qty{53}{L}$ de gas.

\textbf{5.} Unos \qty{53000}{L} $= \qty{53}{m^3}$. La bodega tiene
$4 \times 6 \times 2.5 = \qty{60}{m^3}$: el gas desprendido casi la
llenaría.

\textbf{6.} $\mathrm{C_6H_{12}O_6} + 6\,\mathrm{O_2} \to
6\,\mathrm{CO_2} + 6\,\mathrm{H_2O}$.

\textbf{7.} Oxígeno $6 \times 1.11 = \qty{6.67}{mol}$, $6.67 \times 32
\approx \qty{213}{g}$; dióxido de carbono $6.67 \times 44 \approx
\qty{293}{g}$.

\textbf{8.} Tres veces más: la respiración convierte los seis carbonos
de la glucosa en dióxido de carbono, y la \omterm{def:g10:cell-metabolism:fermentation}{fermentación} solo dos; los
otros cuatro se quedan en el etanol.

\textbf{9.} Con la respiración, quince veces más: cada glucosa da
quince veces más ATP.

\textbf{10.} El etanol solo lo produce la \omterm{def:g10:cell-metabolism:fermentation}{fermentación}; una cuba
aireada daría dióxido de carbono, agua y levadura, y nada de vino.

\textbf{11.} $1.11 \times 70 \approx \qty{78}{kJ}$ por litro;
\qty{78000}{kJ}, es decir \qty{78}{MJ}, para la cuba.

\textbf{12.} $78/4.2 \approx \qty{18.5}{\celsius}$.

\textbf{13.} $20 + 18.5 \approx \qty{38.5}{\celsius}$: al borde del
límite de la levadura, y una bodega cálida lo rebasaría. Una cuba
grande pierde calor despacio porque su superficie es pequeña frente a
su volumen, así que hay que refrigerarla a propósito.

\textbf{14.} Las \omterm{def:g10:cell-metabolism:metabolism}{enzimas} son \omterm{def:g10:chemistry-of-life:families}{proteínas}; el calor las despliega y dejan
de funcionar. Sin sus \omterm{def:g10:cell-metabolism:metabolism}{enzimas}, la levadura no puede transformar el
azúcar, por mucho que quede.

\textbf{15.} En los enlaces químicos del etanol: unos \qty{2800}{kJ}
por mol de glucosa, y por eso el etanol arde.

\textbf{16.} En el suelo, en una capa que engorda a medida que avanza
la \omterm{def:g10:cell-metabolism:fermentation}{fermentación}. Una vela llevada a ras de suelo se apaga allí donde el
dióxido de carbono ha desplazado al oxígeno --- un aviso antes de que
una persona se vea afectada.

\textbf{17.} A la \omterm{def:g10:cell-metabolism:fermentation}{fermentación} láctica. Rinde quince veces menos ATP;
el cerebro y el corazón no pueden cubrir así sus necesidades, y el
ácido láctico se acumula. La pérdida de conciencia llega en uno o dos
minutos.

\textbf{18.} El etanol en concentración alta daña las membranas y las
\omterm{def:g10:cell-metabolism:metabolism}{enzimas} de la levadura: el propio producto de las \omterm{def:g10:cells-common-unit:cell}{células} envenena su
\omterm{def:g10:cell-metabolism:metabolism}{metabolismo}.

\textbf{19.} Las burbujas son dióxido de carbono de la \omterm{def:g10:cell-metabolism:fermentation}{fermentación}. El
etanol se evapora en el horno (hierve a \qty{78}{\celsius}).

\textbf{20.} Alrededor del 13\% de alcohol en volumen y \qty{53}{L} de
gas por litro de mosto; ambos se siguen de que la levadura fermenta,
sin oxígeno, cada glucosa en dos etanoles y dos dióxidos de carbono.
\end{solution}
